\begin{abstract}
本文围绕众擎PM01机器人拳击竞技中的攻击选择、防守匹配、单场决策与故障资源调度四个核心问题，构建了一套以五状态转移为基础的建模与求解框架。首先，给出\textbf{反冲表}（攻击后自身稳定性分布）与\textbf{正冲表}（对手受击后的状态转移分布）的构建思路，据此评估13种攻击动作的综合效用并形成优先序列。其次，基于格挡概率、减伤映射与反击窗口等指标，建立防守匹配模型，为不同来袭攻击生成最优防守响应。再次，在不考虑故障的单场场景下，将攻防决策建模为马尔可夫决策过程（MDP），并通过价值迭代求解不同风险偏好下的最优策略。最后，面向三局两胜（BO3）赛制，综合人工复位、战术暂停与紧急维修等资源约束，建立线性规划与MDP耦合的调度模型。需要说明的是，本文现阶段以方法流程验证为主，结果数据主要用于算法检验，不作为高保真实测结论。
\end{abstract}

\section{引言}
众擎PM01机器人拳击竞技具有强对抗、强时序和强不确定性的特点，机器人需要在攻击收益、防守稳定性与资源消耗之间持续权衡。围绕赛题需求，本文重点研究以下四个问题：
\begin{enumerate}[label=(\arabic*)]
    \item 针对附录一中的13种攻击动作，综合冲击效果与自身稳定性，确定优先采用的攻击动作及其强弱顺序；
    \item 针对附录二中的22种防守动作，为每种攻击匹配最优的防守响应；
    \item 在不考虑故障的单场比赛中，建立最优作战策略模型；
    \item 在BO3赛制下，考虑故障风险与有限恢复资源，给出资源使用的最优决策。
\end{enumerate}
在方法层面，本文以物理机理和状态转移为主线，严格区分\textbf{反冲表}（自身后坐力）与\textbf{正冲表}（对手受击效应），并在此基础上串联攻防评估、序贯决策与资源调度三个层级。受时间与实验条件限制，当前尚未完成完整高保真实测统计；文中 JSON 参数由经验先验与随机化方式构造，主要用于验证建模框架与算法流程的可行性。
